\chapter{Conclusion and Future Work}
\label{chap:conclusion}

%OVERVIEW AND RESULTS
Throughout this dissertation, we analyzed two main strategies used to represent real values computationally: nested intervals and modulus of convergence. As we have seen, the two strategies have different advantages and disadvantages. Modulus of convergence representation is a more direct application of their shared mathematical foundation -- namely, Cauchy sequences --, while nested intervals representation, also based on domain theory, has better performance. We then studied hybridness and hybrid systems in order to build Jaguar, a hybrid simulator that uses exact real arithmetic for its calculations. After exploring its inner workings and functionalities, we tested the various exact real arithmetic libraries, both by themselves and with the hybrid simulator. Of the 5 libraries, only the two using nested intervals representation were viable when running a case study of an inverted pendulum, with \textit{aern2-real} being able to run for long enough to notice discrepancies from a simulation using double precision floating-point numbers.

%FUTURE WORK
Many paths of investigation are still unexplored. Other representation strategies, such as continued fractions and dedekind cuts, have also been suggested as basis for exact real arithmetic. Though these other strategies do not have readily available Haskell libraries at time of writing (which is why they were not included in this dissertation), it would still be interesting to expand the analysis done here to these other methods, both to their theory and practical implementations.

As for the hybrid simulator, Jaguar, there is still plenty of room for improvement. An especially important component to improve is the differential equation solver, considering how crucial it is for the simulator's performance. The most direct way to optimize the solver is to symbolically solve the equations instead of using the Taylor method, completely circumventing the need for automatic differentiation and limits and massively reducing the number of arithmetic operations of each differential statement. Of course, not all differential equations are symbolically solvable, so the current approach could still be used in those cases. Some hybrid programs might also benefit from compiler-like optimizations on all arithmetic expressions (e.g. factoring out and reusing common subexpressions), since exact real arithmetic is highly sensitive to the quantity, order and associativity of operations. Adding this analysis to the parser could be an easy way to improve Jaguar's performance.

Besides optimizations, the simulator could also be improved by adding other features and language constructs, such as Lince's events (which are for statements that run for as long as a condition is true) or some equivalent specific to computable real numbers. Currently, there is also no mechanism to detect and manage arithmetic errors such as division by zero. This problem is especially difficult to manage in exact real arithmetic due to the semidecidable comparison, which makes it impossible to simply check if the divisor is zero. Even worse, evaluating the result of a division by zero might not even terminate, as the division uses smaller and smaller approximations of the divisor, driving the value of the result infinitely high.

Another important area requiring future study is memory usage, which is mostly disregarded in this dissertation but would be of critical importance in an effort to scale up programs and computations. The current technologies of exact real arithmetic would greatly benefit from a memory usage analysis similar to the execution time analysis presented in \cref{sec:package_benches}. The simulator itself should also be optimized for memory efficiency in the future so it can handle more complicated hybrid programs without running out of memory, a sad fate my 8GB laptop often encountered during development and testing.

%CONCLUSION
In conclusion, despite being simple and flexible in application, exact real arithmetic is still too performance heavy, being best employed in specific situations for small but critical calculations. This makes exact real arithmetic, and Jaguar in particular, useful only for formal verifications and theoretical problems, while being too slow for "real world" applications. Nevertheless, this dissertation shows how this technology can be used and hopefully serves as the groundwork for more complex applications in the future.
